\begin{proof}[Proof of \cref{obj:lem:legal-iv-mixture-full-elbow}]
\begin{enumerate}
\item
We first construct a ceiling for amplitudes that gives uniform model membership. Define
\[
c_{\max}
=
\min\left\{
\frac1{100},\,
1-c_f,\,
C_f-1,\,
\frac{L-1}{1+2\pi},\,
\frac{L-3/4}{6\pi}
\right\}.
\]
The hypotheses imply that every entry in this minimum is positive, so
\[
0<c_{\max}\le\frac1{100}<1.
\]
For the membership verification, let
\[
0<c_{\star}\le c_{\max},\qquad
n\ge256,\qquad
0<a\le\frac14,\qquad
\tau\in[-1,1],\qquad
\lambda\in\{-1,1\}^{K_{\mathrm L}}.
\]
Use the construction of \cref{obj:def:legal-iv-mixture}, and write
\[
K_{\mathrm L}=\lceil n^{4/3}\rceil,\qquad
b=\bar a=\max\{a,n^{-1/3}\},\qquad
h=h_n=c_{\star}K_{\mathrm L}^{-1/8},\qquad
u=u_{\lambda},\qquad v=v_{\lambda,\tau}=\tau u.
\]
Since \(n\ge64\),
\[
n^{-1/3}\le64^{-1/3}=\frac14,
\qquad
0<a\le b\le\frac14.
\]
Moreover,
\[
1\le n^{4/3}\le K_{\mathrm L},
\qquad
0<h\le c_{\star}\le\frac1{100},
\]
and monotonicity of the negative power gives
\[
h^2=c_{\star}^2K_{\mathrm L}^{-1/4}
\le c_{\star}^2n^{-1/3}
\le\frac{b}{10000}.
\]
Thus
\[
h^2\le\frac1{10000},
\qquad
\frac14-h^2\ge\frac14-\frac1{10000}\ge\frac1{10000},
\]
and consequently
\[
c_{\star}>0,\qquad
h\le\frac3{16},\qquad
h^2\le\frac{b}{10000}
\le b\left(\frac14-h^2\right).
\]
This proves admissibility throughout the prescribed ranges.

For clarity, the tiled perturbation has the following definition on all of \(\mathbb R\):
\[
u(x)=
\begin{cases}
h\lambda_{\ell}B(K_{\mathrm L}x-\ell),
&x\in\operatorname{cell}(K_{\mathrm L},\ell),
\quad 0\le\ell<K_{\mathrm L},\\
0,&x\notin[0,1],
\end{cases}
\qquad
B(t)=\sin(2\pi t),\quad t\in[0,1].
\]
Here the cells are those of \cref{obj:synth_18}, and \(B\) is the restriction of the bump in \cref{obj:synth_16}. Their disjointness, the endpoint zeros, and the bump normalization give
\[
|v(x)|\le |u(x)|\le h
\quad(x\in\mathbb R),
\qquad
\int_0^1u(x)\,dx
=
\frac{h}{K_{\mathrm L}}
\sum_{\ell=0}^{K_{\mathrm L}-1}
\lambda_{\ell}\int_0^1B(t)\,dt
=0.
\]
In particular, the construction's covariate densities and encouragement probabilities satisfy
\[
f_{\mathrm S}=1,\qquad f_{\mathrm T}=1+u,\qquad
e_z^{(\lambda)}=\frac12+(2z-1)u,\quad z\in\{0,1\},
\]
\[
\int_0^1f_{\mathrm T}(x)\,dx=1,
\qquad
c_f\le1-c_{\star}\le f_{\mathrm T}(x)
\le1+c_{\star}\le C_f,
\]
and
\[
\frac14\le\frac12-h
\le e_z^{(\lambda)}(x)
\le\frac12+h\le\frac34.
\]
The constant source density satisfies the same density envelopes. These are the density and overlap bounds required in
\cref{obj:ass:source-density-bounds,obj:ass:target-density-bounds,obj:ass:propensity-overlap}.
% lean: CausalSmith.Stat.TransportCaceRoughnuisanceLength.exists_membership_amplitude_ceiling
% lean: CausalSmith.Stat.TransportCaceRoughnuisanceLength.lower_admissible_of_small_amplitude
% lean: CausalSmith.Stat.TransportCaceRoughnuisanceLength.legalIVComponent_density_overlap_of_ceiling

\item
We verify the Hölder conditions with the modulus used by the construction. Set
\[
\beta=\frac18.
\]
On a single cell, the sine Lipschitz bound gives
\[
|u(x)-u(x')|
\le 2\pi hK_{\mathrm L}|x-x'|.
\]
The same bound holds across different cells. Indeed, after ordering \(x\le x'\), suppose they belong to cells \(\ell<\ell'\), and define the intervening endpoints by
\[
x_{\ell,+}=\frac{\ell+1}{K_{\mathrm L}},
\qquad
x_{\ell',-}=\frac{\ell'}{K_{\mathrm L}}.
\]
The bump vanishes at these endpoints, and
\[
x\le x_{\ell,+}\le x_{\ell',-}\le x'.
\]
Therefore
\[
\begin{aligned}
|u(x)-u(x')|
&\le |u(x)|+|u(x')|\\
&\le 2\pi hK_{\mathrm L}
\left\{(x_{\ell,+}-x)+(x'-x_{\ell',-})\right\}\\
&\le 2\pi hK_{\mathrm L}|x-x'|.
\end{aligned}
\]
The opposite ordering follows by symmetry.

For \(x,x'\in[0,1]\), define
\[
d_{x,x'}=|x-x'|.
\]
When \(d_{x,x'}=0\), the Hölder bound is immediate. If
\(0<K_{\mathrm L}d_{x,x'}\le1\), then
\[
\begin{aligned}
|u(x)-u(x')|
&\le2\pi c_{\star}K_{\mathrm L}^{-1/8}
(K_{\mathrm L}d_{x,x'})\\
&\le2\pi c_{\star}K_{\mathrm L}^{-1/8}
(K_{\mathrm L}d_{x,x'})^{1/8}
=2\pi c_{\star}d_{x,x'}^{1/8}.
\end{aligned}
\]
If \(K_{\mathrm L}d_{x,x'}>1\), the supremum bound instead gives
\[
|u(x)-u(x')|
\le2h
\le2c_{\star}d_{x,x'}^{1/8}.
\]
Combining the cases yields the uniform modulus
\[
|u(x)-u(x')|
\le c_{\star}(1+2\pi)|x-x'|^{1/8}.
\]
It also proves continuity, including at cell boundaries. Hence
\[
[f_{\mathrm T}]_{\beta},\,[e_1^{(\lambda)}]_{\beta}
\le c_{\star}(1+2\pi)\le L-1<L,
\]
while
\[
\|f_{\mathrm T}\|_{\infty}\le1+c_{\star}\le L,
\qquad
\|e_1^{(\lambda)}\|_{\infty}\le\frac12+c_{\star}\le L.
\]
The source density is constant with supremum norm one and seminorm zero.

The baseline cells of \cref{obj:synth_15} and finite summation in
\cref{obj:def:legal-iv-mixture} give the exact arm means
\[
m_{Dz}(x)=r_z^{(b)}=\frac{1+(2z-1)b}{2},
\qquad
m_{Yz}(x)=
\frac12+\frac{\tau u(x)}{1+2(2z-1)u(x)},
\quad z\in\{0,1\}.
\]
The receipt means are constant and belong to \([3/8,5/8]\). For the outcome means, the denominator satisfies
\[
1+2(2z-1)u(x)\ge1-\frac2{100}=\frac{49}{50}.
\]
Subtracting the two rational expressions gives
\[
m_{Yz}(x)-m_{Yz}(x')
=
\frac{\tau\{u(x)-u(x')\}}
{\{1+2(2z-1)u(x)\}\{1+2(2z-1)u(x')\}}.
\]
Since \(|\tau|\le1\) and \((49/50)^2\ge1/2\),
\[
|m_{Yz}(x)-m_{Yz}(x')|
\le2|u(x)-u(x')|.
\]
Taking \(u(x')=0\) in the same rational comparison gives
\[
\left|m_{Yz}(x)-\frac12\right|
\le2|u(x)|\le\frac1{50},
\]
so the outcome means belong to \([12/25,13/25]\). Their Hölder seminorms satisfy
\[
[m_{Yz}]_{\beta}
\le2c_{\star}(1+2\pi)
\le6\pi c_{\star}
\le L-\frac34<L.
\]
Their supremum norms are at most \(13/25<L\). Thus all density, propensity, and arm-mean Hölder conditions hold, with the range conditions in
\cref{obj:ass:arm-means-holder}. The corresponding center functions are the constants
\[
f_{\mathrm S}=f_{\mathrm T}=1,\qquad
e=\frac12,\qquad
m_{Dz}=r_z^{(b)},\qquad
m_{Yz}=\frac12,
\]
and satisfy the same conditions.
% lean: CausalSmith.Stat.TransportCaceRoughnuisanceLength.tiledPerturbation_holder_modulus
% lean: CausalSmith.Stat.TransportCaceRoughnuisanceLength.rational_arm_mean_abs_sub_le
% lean: CausalSmith.Stat.TransportCaceRoughnuisanceLength.legalIVComponent_arm_means_holder

\item
We next verify the primitive probability masses and calculate their complier margins. By admissibility and \(|v|\le|u|\),
\[
|u(x)v(x)|\le u(x)^2\le h^2
\le b\left(\frac14-h^2\right)
\le b\left(\frac14-u(x)^2\right),
\]
and
\[
\frac14-u(x)^2\ge\frac14-h^2>0.
\]
It follows that
\[
\left|
\frac{u(x)v(x)}
{2\{1/4-u(x)^2\}}
\right|\le\frac b2.
\]
Consequently the complier margins \(M_{dy}(x)\) defined in
\cref{obj:def:legal-iv-mixture} satisfy
\[
M_{dy}(x)\ge0,
\qquad
M_{d0}(x)+M_{d1}(x)=b,
\quad d,y\in\{0,1\}.
\]

For always-takers and never-takers, the relevant masses in that construction are
\[
\frac{1-b}{4}
+\frac{v(x)(2y-1)}{4e_0^{(\lambda)}(x)},
\qquad
\frac{1-b}{4}
+\frac{v(x)(2y-1)}{4e_1^{(\lambda)}(x)},
\quad y\in\{0,1\}.
\]
Using \(|u|,|v|\le3/16\), we obtain
\[
e_z^{(\lambda)}(x)\ge\frac5{16},
\qquad
\left|\frac{v(x)}{4e_z^{(\lambda)}(x)}\right|
\le\frac{3/16}{4(5/16)}=\frac3{20},
\qquad
\frac{1-b}{4}\ge\frac3{16}.
\]
Thus these masses are nonnegative. Summing over \(y\), the perturbations cancel, and each of the two strata has total mass \((1-b)/2\). The independent complier coupling has joint masses
\[
P(C=1,Y(0)=y_0,Y(1)=y_1\mid X=x)
=
\frac{M_{0y_0}(x)M_{1y_1}(x)}{b},
\quad y_0,y_1\in\{0,1\},
\]
whose sum is
\[
\frac{
\left(\sum_{y_0=0}^1M_{0y_0}(x)\right)
\left(\sum_{y_1=0}^1M_{1y_1}(x)\right)
}{b}
=b.
\]
The three principal strata therefore have total mass
\[
\frac{1-b}{2}+b+\frac{1-b}{2}=1.
\]
The remaining primitive cells have mass zero, as specified by the construction.

The receipt assignments \((1,1)\), \((0,0)\), and \((0,1)\) give monotonicity. Conditional independent encouragement, followed by
\[
D=D(Z),\qquad Y=Y(D),
\]
gives receipt consistency, outcome consistency, instrument randomization, and exclusion. All potential outcomes are binary. These verify
\cref{obj:ass:receipt-consistency,obj:ass:outcome-consistency,obj:ass:instrument-randomization,obj:ass:exclusion,obj:ass:monotonicity}.

Finite summation of the same complier masses gives, for either population and every \(x\in[0,1]\),
\[
\mathbb E[C\mid X=x,S=s]=b,
\qquad s\in\{0,1\},
\]
and
\[
\begin{aligned}
\mathbb E[(Y(1)-Y(0))C\mid X=x,S=s]
&=
\sum_{y_0,y_1\in\{0,1\}}
(y_1-y_0)\frac{M_{0y_0}(x)M_{1y_1}(x)}b\\
&=M_{11}(x)-M_{01}(x)
=-\frac{u(x)v(x)}{1/4-u(x)^2}.
\end{aligned}
\]
These are also versions of the conditional means. Indeed, the construction uses the same primitive conditional mass function in both populations. Define their covariate densities by
\[
f_{S=0}(x)=1+u(x),\qquad f_{S=1}(x)=1,
\quad x\in[0,1].
\]
Integrating the finite sums over any Borel set \(E_X\subseteq\mathbb R\) yields
\[
\int_{\{X\in E_X\}} C\,dQ_{\lambda,\tau}(\,\cdot\mid S=s)
=
\int_{E_X\cap[0,1]}b f_{S=s}(x)\,dx
\]
and
\[
\int_{\{X\in E_X\}}(Y(1)-Y(0))C\,
dQ_{\lambda,\tau}(\,\cdot\mid S=s)
=
\int_{E_X\cap[0,1]}
-\frac{u(x)v(x)}{1/4-u(x)^2}f_{S=s}(x)\,dx.
\]
The displayed conditional-mean functions are measurable and bounded. These identities verify their conditional-mean property and both transport requirements.

The arm-mean formulas from the preceding step give
\[
\Delta_D(x)=b,\qquad
\Delta_Y(x)=-\frac{u(x)v(x)}{1/4-u(x)^2}.
\]
Hence
\[
\mu(Q_{\lambda,\tau})
=
\int_0^1(1+u(x))b\,dx=b.
\]
The target complier share is also \(b\), by the conditional-mean identity with \(E_X=\mathbb R\).

Define the separation functional by
\[
J_n
=
h^2\int_0^1
\frac{B(t)^2}{1/4-h^2B(t)^2}\,dt.
\]
The reflection identity \(B(1-t)=-B(t)\) from \cref{obj:synth_16} implies
\[
\int_0^1
\frac{B(t)^3}{1/4-h^2B(t)^2}\,dt
=
-\int_0^1
\frac{B(t)^3}{1/4-h^2B(t)^2}\,dt
=0.
\]
Changing variables separately on each cell therefore gives
\[
\int_0^1\frac{u(x)^2}{1/4-u(x)^2}\,dx=J_n,
\qquad
\int_0^1\frac{u(x)^3}{1/4-u(x)^2}\,dx=0.
\]
Since \(v=\tau u\), the target complier-weighted outcome effect equals
\[
\begin{aligned}
\mathbb E_{Q_{\lambda,\tau}}[(Y(1)-Y(0))C\mid S=0]
&=-\tau\int_0^1
\left\{
\frac{u(x)^2}{1/4-u(x)^2}
+\frac{u(x)^3}{1/4-u(x)^2}
\right\}\,dx\\
&=-\tau J_n.
\end{aligned}
\]
Dividing by the positive target complier share gives
\[
\theta(Q_{\lambda,\tau})=-\frac{\tau J_n}{b}.
\]
In particular,
\[
J_n=-b\,\theta(Q_{\lambda,1}).
\]
For the center, the same calculations with zero perturbations give
\[
\mu(P_{\star})=b,\qquad \theta(P_{\star})=0.
\]
% lean: CausalSmith.Stat.TransportCaceRoughnuisanceLength.primitiveMass_nonneg_of_bounds
% lean: CausalSmith.Stat.TransportCaceRoughnuisanceLength.lowerPointwiseMean_complier_margins
% lean: CausalSmith.Stat.TransportCaceRoughnuisanceLength.legal_mixture_margins

\item
Independent source and target product sampling supplies
\cref{obj:ass:source-iid,obj:ass:target-iid,obj:ass:sample-independence}.
The preceding calculations establish the density identities, support, regularity, causal restrictions, conditional-mean transport, and positive strength required by
\cref{obj:def:model-class}. Thus
\[
Q_{\lambda,\tau}\in\mathcal M_n,\qquad
P_{\star}\in\mathcal M_n.
\]
Together with
\[
a\le b=\mu(Q_{\lambda,\tau})=\mu(P_{\star})\le\frac14,
\]
this proves, by \cref{obj:def:strength-slice},
\[
Q_{\lambda,\tau}\in\mathcal M_n(a),
\qquad
P_{\star}\in\mathcal M_n(a).
\]
We have established these memberships uniformly for every
\(c_{\star}\in(0,c_{\max}]\), every admissible \(n,a\), and every \(\lambda,\tau\).
% lean: CausalSmith.Stat.TransportCaceRoughnuisanceLength.legal_mixture_membership

\item
We now choose one amplitude within this membership ceiling. Since \(\alpha<1\),
\[
\exp\{C_{\mathrm{mix}}\,0^4\}=1<1+(1-\alpha)^2,
\qquad
C_{\mathrm{mix}}=\frac{19321}{1800}.
\]
Continuity at zero supplies \(\delta>0\) such that
\[
0\le c_{\mathrm{test}}<\delta
\quad\Longrightarrow\quad
\exp\{C_{\mathrm{mix}}c_{\mathrm{test}}^4\}
<1+(1-\alpha)^2.
\]
Choose
\[
c_{\star}
=
\min\left\{c_{\max},\,\frac1{100},\,\frac{\delta}{2}\right\}.
\]
Then
\[
0<c_{\star}\le c_{\max}<1,\qquad
c_{\star}\le\frac1{100},\qquad
\exp\{C_{\mathrm{mix}}c_{\star}^4\}\le1+(1-\alpha)^2.
\]
This choice depends on \(\alpha,c_f,C_f,L\) and is uniform in \(n,a,\lambda,\tau\). The admissibility calculation in the first step applies to this amplitude for every prescribed \(n,a\).
% lean: CausalSmith.Stat.TransportCaceRoughnuisanceLength.calibrated_mixture_amplitude

\item
Fix \(n\ge256\) and \(0<a\le1/4\). We bound the common separation functional. The ceiling bounds and \(n^{4/3}\ge1\) give
\[
n^{4/3}\le K_{\mathrm L}
\le n^{4/3}+1\le2n^{4/3}.
\]
Since \(h^2\le1/10000\le1/8\) and \(B(t)^2\le1\),
\[
\frac18\le\frac14-h^2B(t)^2\le\frac14,
\]
so
\[
4B(t)^2
\le\frac{B(t)^2}{1/4-h^2B(t)^2}
\le8B(t)^2.
\]
Using \(\int_0^1B(t)^2\,dt=1/2\), supplied by \cref{obj:synth_16}, we obtain
\[
2
\le\int_0^1\frac{B(t)^2}{1/4-h^2B(t)^2}\,dt
\le4,
\qquad
2h^2\le J_n\le4h^2.
\]
Finally,
\[
(2n^{4/3})^{-1/4}
\le K_{\mathrm L}^{-1/4}\le n^{-1/3},
\]
and therefore
\[
2^{3/4}c_{\star}^2n^{-1/3}
=
2c_{\star}^2(2n^{4/3})^{-1/4}
\le J_n
\le4c_{\star}^2n^{-1/3}.
\]
% lean: CausalSmith.Stat.TransportCaceRoughnuisanceLength.legal_mixture_separation

\item
Fix \(\tau\in[-1,1]\). We calculate the distance of the uniform sign mixture from the center. With respect to Lebesgue measure on \([0,1]\) and counting measure on the three binary source coordinates, the center source density is
\[
p_{\star,\mathrm S}(x,z,d,y)
=\frac12R_z^{(b)}(d,y),
\qquad
x\in[0,1],\quad z,d,y\in\{0,1\}.
\]
The baseline formulas in \cref{obj:synth_15} give
\[
R_z^{(b)}(d,y)\ge\frac{1-b}{4}\ge\frac3{16}>0.
\]
The source cells of \cref{obj:def:legal-iv-mixture} consequently yield the one-record likelihood ratios
\[
\operatorname{LR}_{\mathrm S,\lambda,\tau}(x,z,d,y)
=
1+
u_{\lambda}(x)
\frac{(2z-1)R_z^{(b)}(d,y)+\tau(2y-1)/4}
{R_z^{(b)}(d,y)/2},
\]
\[
\operatorname{LR}_{\mathrm T,\lambda}(x)=1+u_{\lambda}(x).
\]
The first formula is used on the center source support
\([0,1]\times\{0,1\}^3\), and the second on the center target support \([0,1]\); take both ratios to be zero off their respective supports. They are nonnegative, measurable, and bounded, and satisfy
\[
Q_{\lambda,\tau,\mathrm S}
=
\operatorname{LR}_{\mathrm S,\lambda,\tau}\,P_{\star,\mathrm S},
\qquad
Q_{\lambda,\tau,\mathrm T}
=
\operatorname{LR}_{\mathrm T,\lambda}\,P_{\star,\mathrm T}.
\]

Define, for \(b\in(0,1/4]\) and \(\tau\in[-1,1]\),
\[
\gamma(b,\tau)
=
\sum_{z,d,y\in\{0,1\}}
\frac{\{(2z-1)R_z^{(b)}(d,y)+\tau(2y-1)/4\}^2}
{R_z^{(b)}(d,y)/2}.
\]
Expanding the square, its three finite sums are
\[
\sum_{z,d,y}
\frac{\{(2z-1)R_z^{(b)}(d,y)\}^2}{R_z^{(b)}(d,y)/2}
=4,
\]
\[
\sum_{z,d,y}
\frac{(2z-1)R_z^{(b)}(d,y)(2y-1)/4}
{R_z^{(b)}(d,y)/2}
=0,
\]
and
\[
\sum_{z,d,y}
\frac{\{(2y-1)/4\}^2}{R_z^{(b)}(d,y)/2}
=
\frac2{1+b}+\frac2{1-b}
=\frac4{1-b^2}.
\]
The middle sum cancels over \(y\); the last uses the four baseline cells of each mass \((1+b)/4\) and \((1-b)/4\). Thus
\[
\gamma(b,\tau)=4+\frac{4\tau^2}{1-b^2}.
\]
Since \(1-b^2\ge15/16\) and \(\tau^2\le1\),
\[
0\le\gamma(b,\tau),
\qquad
\gamma(b,\tau)+1
\le5+\frac{64}{15}=\frac{139}{15}.
\]

For two sign vectors, define their overlap by
\[
\mathscr O(\lambda,\lambda')
=
\sum_{\ell=0}^{K_{\mathrm L}-1}\lambda_{\ell}\lambda'_{\ell},
\qquad
\lambda,\lambda'\in\{-1,1\}^{K_{\mathrm L}}.
\]
Cellwise integration gives
\[
\begin{aligned}
\int_0^1u_{\lambda}(x)u_{\lambda'}(x)\,dx
&=
\frac{h^2}{K_{\mathrm L}}
\sum_{\ell=0}^{K_{\mathrm L}-1}
\lambda_{\ell}\lambda'_{\ell}\int_0^1B(t)^2\,dt\\
&=\frac{h^2}{2K_{\mathrm L}}\mathscr O(\lambda,\lambda').
\end{aligned}
\]
Expanding the likelihood products, the zero-integral linear terms cancel. Consequently,
\[
\int
\operatorname{LR}_{\mathrm S,\lambda,\tau}
\operatorname{LR}_{\mathrm S,\lambda',\tau}\,dP_{\star,\mathrm S}
=
1+\frac{\gamma(b,\tau)h^2}{2K_{\mathrm L}}
\mathscr O(\lambda,\lambda'),
\]
\[
\int
\operatorname{LR}_{\mathrm T,\lambda}
\operatorname{LR}_{\mathrm T,\lambda'}\,dP_{\star,\mathrm T}
=
1+\frac{h^2}{2K_{\mathrm L}}
\mathscr O(\lambda,\lambda').
\]
Both expressions are nonnegative, being integrals of products of nonnegative likelihoods.

For an observed sample
\[
\omega=
\bigl((x_i,z_i,d_i,y_i)_{i=1}^n,\,
(x_r^{\mathrm T})_{r=1}^n\bigr),
\]
define the product and mixture likelihoods by
\[
\operatorname{LR}_{\lambda,\tau}^{(n,n)}(\omega)
=
\prod_{i=1}^n
\operatorname{LR}_{\mathrm S,\lambda,\tau}(x_i,z_i,d_i,y_i)
\prod_{r=1}^n
\operatorname{LR}_{\mathrm T,\lambda}(x_r^{\mathrm T}),
\]
\[
\operatorname{LR}_{\tau}^{\mathrm{mix}}(\omega)
=
2^{-K_{\mathrm L}}
\sum_{\lambda\in\{-1,1\}^{K_{\mathrm L}}}
\operatorname{LR}_{\lambda,\tau}^{(n,n)}(\omega).
\]
Source--target independence and independent sampling within each channel give
\[
M_{\tau}
=
\operatorname{LR}_{\tau}^{\mathrm{mix}}\,P_{\star}^{(n,n)}.
\]
The bounded one-record likelihoods give integrable product pairs. The finite mixture then has an integrable square, since
\[
\bigl(\operatorname{LR}_{\tau}^{\mathrm{mix}}\bigr)^2
=
2^{-2K_{\mathrm L}}
\sum_{\lambda,\lambda'}
\operatorname{LR}_{\lambda,\tau}^{(n,n)}
\operatorname{LR}_{\lambda',\tau}^{(n,n)}.
\]
Its mean under the center is one.

Let \(\lambda,\lambda'\) be independent uniform sign vectors. Expanding this square and using the two channel inner products gives
\[
\begin{aligned}
1+\chi^2(M_{\tau},P_{\star}^{(n,n)})
&=
\mathbb E_{\lambda,\lambda'}
\left(1+
\frac{\gamma(b,\tau)h^2}{2K_{\mathrm L}}
\mathscr O(\lambda,\lambda')\right)^n
\left(1+
\frac{h^2}{2K_{\mathrm L}}
\mathscr O(\lambda,\lambda')\right)^n\\
&\le
\mathbb E_{\lambda,\lambda'}
\exp\left\{
\frac{nh^2\{\gamma(b,\tau)+1\}}
{2K_{\mathrm L}}\mathscr O(\lambda,\lambda')
\right\}.
\end{aligned}
\]
Here \(1+t\le e^t\) is applied to each nonnegative base, followed by multiplication of their \(n\)-th powers.

For each fixed \(\lambda\), the products \(\lambda_{\ell}\lambda'_{\ell}\) are independent uniform signs. Factoring their exponential moment and using the standard inequality
\[
\cosh(t)\le e^{t^2/2},\qquad t\in\mathbb R,
\]
gives
\[
\begin{aligned}
\mathbb E_{\lambda'}
\exp\left\{
\frac{nh^2\{\gamma(b,\tau)+1\}}
{2K_{\mathrm L}}\mathscr O(\lambda,\lambda')
\right\}
&=
\cosh\left(
\frac{nh^2\{\gamma(b,\tau)+1\}}{2K_{\mathrm L}}
\right)^{K_{\mathrm L}}\\
&\le
\exp\left\{
\frac{n^2h^4\{\gamma(b,\tau)+1\}^2}
{8K_{\mathrm L}}
\right\}.
\end{aligned}
\]
Averaging over \(\lambda\) preserves this bound. Finally,
\[
n^2\le K_{\mathrm L}^{3/2},
\qquad
h^4=c_{\star}^4K_{\mathrm L}^{-1/2},
\]
so
\[
\frac{n^2h^4\{\gamma(b,\tau)+1\}^2}{8K_{\mathrm L}}
=
\frac{n^2}{K_{\mathrm L}^{3/2}}
\frac{\{\gamma(b,\tau)+1\}^2}{8}c_{\star}^4
\le
\frac{(139/15)^2}{8}c_{\star}^4
=
\frac{19321}{1800}c_{\star}^4.
\]
This proves
\[
1+\chi^2(M_{\tau},P_{\star}^{(n,n)})
\le\exp\{C_{\mathrm{mix}}c_{\star}^4\}.
\]

With the distance normalization in \cref{obj:synth_17}, absolute continuity and Cauchy--Schwarz yield
\[
\begin{aligned}
\operatorname{TV}(M_{\tau},P_{\star}^{(n,n)})
&=
\frac12\int
\left|\operatorname{LR}_{\tau}^{\mathrm{mix}}-1\right|
\,dP_{\star}^{(n,n)}\\
&\le
\frac12\left\{
\int
\left(\operatorname{LR}_{\tau}^{\mathrm{mix}}-1\right)^2
\,dP_{\star}^{(n,n)}
\right\}^{1/2}\\
&=\frac12\sqrt{\chi^2(M_{\tau},P_{\star}^{(n,n)})}.
\end{aligned}
\]
The amplitude calibration therefore gives
\[
\chi^2(M_{\tau},P_{\star}^{(n,n)})
\le\exp\{C_{\mathrm{mix}}c_{\star}^4\}-1
\le(1-\alpha)^2,
\]
and, because \(1-\alpha>0\),
\[
\operatorname{TV}(M_{\tau},P_{\star}^{(n,n)})
\le\frac{1-\alpha}{2}.
\]
Every estimate is uniform over \(\tau\in[-1,1]\).
% lean: CausalSmith.Stat.TransportCaceRoughnuisanceLength.sourceLikelihood_pair_integral
% lean: Causalean.Stat.Minimax.Mixture.TwoChannel.one_add_chiSqDiv_uniformMixture_twoChannel_sign_le_exp
% lean: CausalSmith.Stat.TransportCaceRoughnuisanceLength.legal_mixture_distance

\item
The likelihood representation above also establishes
\[
M_{\tau}\ll P_{\star}^{(n,n)}.
\]
The finite sum of integrable product pairs proves integrability of
\((\operatorname{LR}_{\tau}^{\mathrm{mix}})^2\); its mean is one, and hence
\[
\int
\left(\operatorname{LR}_{\tau}^{\mathrm{mix}}-1\right)^2
\,dP_{\star}^{(n,n)}
<\infty.
\]
Thus the likelihood ratio minus one is square-integrable.

Finally, \(n\ge256\) gives \(n>0\), and the chosen amplitude satisfies both admissibility and \(c_{\star}<1\). The triple in
\cref{obj:def:legal-iv-mixture} is therefore defined throughout the stated parameter range. Uniform membership, the center margins, and the separation bounds were established above. For each component and each population, the finite-mass calculations give the asserted pointwise averages, while the integral identities establish their conditional-mean versions. Together with the distance and likelihood conclusions, these prove all the claims.
% lean: CausalSmith.Stat.TransportCaceRoughnuisanceLength.legal_mixture_finite_chiSquare
% lean: CausalSmith.Stat.TransportCaceRoughnuisanceLength.legal_iv_mixture_full_elbow
\end{enumerate}
\end{proof}
